% page 632 du fac-simile (image p-631 du scan)
% bloc 152.4 x 229.0 mm ; marge gauche 8.0 mm
\pgc{-12.700mm}{0.000mm}{
\l{\marge[77.393mm]{18.261mm}{\ornamento{12.737mm}{ornements/letroj/litero-X.png}}\marge[76.432mm]{202.097mm}{\ornamento{12.256mm}{ornements/letroj/litero-Y.png}}\cel{3}624}
\l{}
\l{}
\l{}
\l{}
\l{}
\l{}
\l{}
\l{\sou{xantino}.\cel{2}(\sou{kemio})\cel{2}La\cel{2}C5 H4 N4 O2\cel{2}quan kontenas multa}
\l{\cel{3}liquidi de animali : sango, urino, ed anke teo. -}
\l{}
\l{\sou{xantofilo}.\cel{2}(\sou{kemio)}.\cel{2}C40 H56 O2, kristali flava qui fuzas}
\l{\cel{3}ye 172o C., e renkontresas en la klorofilo di la veje-}
\l{\cel{3}tanto-folii. -}
\l{}
\l{\sou{xenio}. I. (\sou{versif.}) Epigramo satira (enduktita da Goethe}
\l{\cel{3}e Schiller. - II. (\sou{biol.}) Influo eventuala da la elemen-}
\l{\cel{3}to maskula, en ula kazi, ad ula parti di la femino qui,}
\l{\cel{3}pro lo, aquiras la karakteri di la maskulo fekundiginta.}
\l{\cel{3}(Dicesas precipue pri la planti di qui la fekundeso de-}
\l{\cel{3}venas de poleno stranjera). -}
\l{}
\l{\sou{xenono}. (\sou{kemio)}. Un ek la gasi skarsa di atmosfero, mono-}
\l{\cel{3}atoma ; atom-pezo 130,2.\cel{2}-\cel{1}\sou{Simbolo kemiala}\cel{1}: X.}
\l{}
\l{\sou{xerofito}.\cel{2}(\sou{biol.})\cel{2}Singla de la vejetanti quin karakter-}
\l{\cel{3}izas la hardeso di la organi diversa e la sekreco sama-}
\l{\cel{3}tampa di olei eteroza abundanta. -}
\l{}
\l{\sou{xifoida}. (\sou{anat.})\cel{2}Dicesas pri la apendico - prolonguro kar-}
\l{\cel{3}tilaga, qua finas la extremajo infra di sternumo. -}
\l{}
\l{\sou{xilofago}. Singla de la insekti qui rodas ligno. - EFIS.}
\l{}
\l{\sou{xilofono}. (\sou{muziko)}.\cel{2}Instrumento ek ligno-tabuleti ne-inter-}
\l{\cel{3}egale longa, subtenata da du suportili, e sur qui onu}
\l{\cel{3}frapas per du vergeti ligna. - DEFIS.}
\l{}
\l{\sou{xilografar}. (\sou{trans.)}\cel{3}Imprimar ek tipi ligna, od ek tabu-}
\l{\cel{3}ldti ligna sur qui esas grabita vorti e figuri. - DEFIRS.}
\l{}
\l{}
\l{}
\l{}
\l{}
\l{}
\l{}
\l{}
\l{}
\l{}
\l{}
\l{ya.\cel{2}(\sou{adverbo)}.\cel{2}Partikulo qua nulkaze uzesas sole, ma kun}
\l{\cel{3}irga dico, e konfirmas ica emfazoze. -\cel{2}D.}
}
